% sections/03-bypass-auth.tex
\section{Sobrepaso del esquema de autenticación con SQL Injection pasiva}

Para esta parte de la práctica se dispuso del formulario de inicio de sesión de la aplicación vulnerable de demostración bancaria publicada en \url{http://testfire.net/bank/login.aspx}, sobre el cual, haciendo uso de técnicas de inyección SQL manual (pasiva, sin herramienta automatizada), se buscó sobrepasar el sistema de autenticación para acceder a la aplicación, teniendo presente que el objetivo no es identificar usuarios por defecto sino hacer uso de la inyección de sentencias SQL.

Sobre el campo de usuario se inyectó la siguiente sentencia, dejando en el campo de contraseña un valor cualquiera (la contraseña queda fuera de la consulta al ir comentada):

\begin{lstlisting}[language=SQL]
' OR '1'='1' --
\end{lstlisting}

\captura{bypass-login.png}{Formulario de inicio de sesión con la sentencia SQL inyectada en el campo usuario}

\captura{bypass-autenticado.png}{Sesión iniciada tras sobrepasar el esquema de autenticación}

\begin{analisis}
\end{analisis}

\pregunta{Realice un análisis de la sentencia SQL usada: ¿por qué estamos sobrepasando el esquema de autenticación de la aplicación?}
\begin{respuesta}
\end{respuesta}

\porque{base del análisis de la sentencia}{%
  La aplicación construye algo equivalente a\\
  \texttt{SELECT * FROM users WHERE username='}$_{input1}$\texttt{' AND
  password='}$_{input2}$\texttt{'}; el payload cierra el literal de texto con
  la primera comilla, inyecta la condición \texttt{OR '1'='1'} (una tautología:
  siempre verdadera) y con \texttt{--} comenta el resto de la línea original,
  eliminando la comilla final y la verificación de la contraseña. La cláusula
  \texttt{WHERE} queda así siempre verdadera, la consulta devuelve el primer
  usuario de la tabla y la aplicación lo trata como sesión válida. Redáctalo
  con esa cadena: concatenación sin parametrización, tautología y comentario;
  y cierra con que el control ausente es validar la identidad contra el
  resultado de una consulta consultada por credenciales, no confiar en que la
  consulta devuelva vacío.%
}
